%%%%%%%%%%%%%%%%%%%%%%%%%%%%
%%%%%     PREAMBLE     %%%%%
%%%%%%%%%%%%%%%%%%%%%%%%%%%%

\documentclass[12pt]{article}

%Define Margins using Geometry Package
\usepackage[top=1.0in,
bottom=1.0in,
right=1.0in,
left = 1.0in]{geometry}


%math
\usepackage{mathtools}


%Set up links and internal references
\usepackage[colorlinks=true,linkcolor=blue, urlcolor=blue]{hyperref} %hide the ugly red boxes
\usepackage[noabbrev]{cleveref} %don't abbreviate ``figure'', etc.
\usepackage{url} %for urls

%Remove natural indentation for itemize and enumerate environments
\usepackage{enumitem} %package to control itemize/enumerate behavior
\setlist[itemize]{leftmargin=*} %for itemize
\setlist[enumerate]{leftmargin=*} %for enumerate
%Swap out the itemize symbol for an N-dash rather than a bullet (does not require enumitem package)
\def\labelitemi{--}

\usepackage{multicol} %Allow for two columns in the middle of a paper. 


\usepackage{titling}
\setlength{\droptitle}{-2.5cm}

%Set up fancy header and footer
\usepackage{fancyhdr}
\pagestyle{fancy}
%Fancy Header Content
\fancyfoot[L]{ME EN 497R}
\fancyfoot[C]{}
\fancyfoot[R]{Rotor Design - Optimization: Part V}
%Fancy Footer Content
% \fancyfoot[L,C]{} %make bottom left and center empty
% \fancyfoot[R]{\thepage}
\title{Rotor Design - Optimization: Part V}
\date{}

%Center section headings
\usepackage{sectsty}
\subsectionfont{\centering}
%Remove section numbering
\setcounter{secnumdepth}{0}




%%%%%%%%%%%%%%%%%%%%%%%%%%%%%%%%
%%%%%     END PREAMBLE     %%%%%
%%%%%%%%%%%%%%%%%%%%%%%%%%%%%%%%


%BEGIN Document Environment
\begin{document}
\maketitle
\vspace{-3.5cm}

This assignment is comprised of three parts, the first of which is to produce a complete draft and submit it for formal review by graduate students in the FLOW lab other than your assigned mentor. Finally you will present your work.
	
\section{Part I: Draft and Review}

 Complete the following steps:


%%%%%%%%%%%%%%%%%%%%%%%%%%%%%%%%%%%%%
%%%%%      Research Skills      %%%%%
%%%%%%%%%%%%%%%%%%%%%%%%%%%%%%%%%%%%%

\begin{enumerate}[label=\alph*.]
	\item Ask 2 graduate students in the FLOW Lab, who are not your graduate student mentor, to review your penultimate draft of your final report.
	\item Write a brief document (bullets are good) compiling the feedback as well as describing your actions to address the feedback (or why you didn't).
	\item Create a ``Diff'' document comparing your penultimate and final drafts of your final report.
\end{enumerate}


%%%%%%%%%%%%%%%%%%%%%%%%%%%%%%%%%%%%%%
%%%%%      Final Submission      %%%%%
%%%%%%%%%%%%%%%%%%%%%%%%%%%%%%%%%%%%%%

\section{Part II: Final Submission}

Your final submission should be a complete and polished final draft of your semester final technical report.
It should include all the improvements made up to this point as well as improvements based on the feedback you received in the first part of this assignment.
You will be graded on completeness and quality of your final report.


\section{Part III: Presentation}
Your final is a presentation of your work to the lab (minimally all other students doing the 497R, their graduate student mentors, and Dr. Ning will attend). Your presentation should be 5-7 minutes long. Guidelines and tips for developing a presentation can be found in the \href{https://flow.byu.edu/onboarding/presentations/}{FLOW Lab Onboarding Guide}. Your presentation should focus on your methods and results; it should not be a summary of the semester (if you find yourself saying "then I did this", then you are probably off the mark). You will fill out a when-is-good to select a time during finals week to present your work. 
	
\end{document}